\chapter{Βιβλιογραφική Ανασκόπηση}
\label{ch:literature}

\section{Εισαγωγή}

Η βραχυπρόθεσμη πρόβλεψη τιμών ηλεκτρικής ενέργειας (\textit{Electricity Price Forecasting}, EPF)
αποτελεί ενεργό πεδίο έρευνας εδώ και δύο δεκαετίες~\cite{weron2014electricity}. Ο Weron
(2014)~\cite{weron2014electricity} παρέχει μία εκτεταμένη επισκόπηση, οργανώνοντας τα μοντέλα σε
πέντε κατηγορίες: multi-agent, fundamental, reduced-form, statistical και computational intelligence.
Η πιο πρόσφατη ανασκόπηση του Lago et al.\ (2021)~\cite{lago2021} επικεντρώνεται στους αλγόριθμους
τελευταίας γενιάς και προτείνει ένα open-access benchmark για αντικειμενική σύγκριση.

Η πρόβλεψη φορτίου (\textit{Electricity Load Forecasting}, ELF) έχει μακρύτερη ιστορία~\cite{gross1987},
με τα πρώτα στατιστικά μοντέλα να εμφανίζονται στα τέλη του 1980, και η εφαρμογή τεχνικών ML να
παρουσιάζει ταχεία ανάπτυξη από τα μέσα του 2000~\cite{hong2020}.


\section{Παραδοσιακά Στατιστικά Μοντέλα}
\label{sec:stat_models}

\subsection{ARIMA/SARIMA και παραλλαγές}

Τα μοντέλα ARIMA (Autoregressive Integrated Moving Average) και η εποχιακή επέκτασή τους SARIMA
αποτελούν κλασικά εργαλεία για χρονοσειρές~\cite{box2015}. Στο πλαίσιο EPF, λόγω της ισχυρής
ωριαίας και ημερήσιας εποχιακότητας, η επιλογή κατάλληλης δομής $(p,d,q)(P,D,Q)_{24/168}$ είναι
κρίσιμη. Ο Kanousis (2025)~\cite{kanousis2025} εφαρμόζει SARIMAX για την ελληνική αγορά με όρους
Fourier αντί ημερολογιακών μεταβλητών, επιτυγχάνοντας ~19\% βελτίωση έναντι της persistence μεθόδου.

Παρόλα αυτά, τα ARIMA μοντέλα έχουν δύο σοβαρούς περιορισμούς: δεν μπορούν να συλλάβουν μη-γραμμικές
σχέσεις και δυσκολεύονται να ενσωματώσουν πολλαπλές εξωγενείς μεταβλητές αποτελεσματικά.

\subsection{GARCH και Υβριδικά Στατιστικά Μοντέλα}

Τα μοντέλα GARCH (Generalized Autoregressive Conditional Heteroskedasticity) χρησιμοποιούνται κυρίως
για τη μοντελοποίηση της μεταβλητότητας (volatility) των τιμών~\cite{weron2014electricity}. Στην
πράξη, συνδυάζονται με ARIMA (ARMA-GARCH) για ταυτόχρονη εκτίμηση μέσης τιμής και διακύμανσης,
χρήσιμο για διαχείριση κινδύνου.


\section{Μοντέλα Μηχανικής Μάθησης}
\label{sec:ml_models}

\subsection{Δέντρα Αποφάσεων και Μέθοδοι Boosting}

Τα μοντέλα βασισμένα σε δέντρα αποφάσεων, και ιδίως οι συναθροιστικές μέθοδοι
(ensemble methods), έχουν κυριαρχήσει στο EPF benchmark~\cite{lago2021}. Χαρακτηρίζονται από ισχυρή
ικανότητα σύλληψης μη-γραμμικών σχέσεων χωρίς εκτεταμένη προεπεξεργασία.

\textbf{Random Forest (RF):} Εισήχθη από τον Breiman (2001)~\cite{breiman2001rf} και αποτελεί
bagging ensemble δέντρων απόφασης. Αποδεικνύει καλή ανθεκτικότητα σε θόρυβο και υπεραρμογή,
ιδιαίτερα με \textit{subsampling} χαρακτηριστικών ανά κόμβο.

\textbf{Gradient Boosting (XGBoost, LightGBM):} To XGBoost~\cite{chen2016xgboost} αξιοποιεί
δεύτερης τάξης κλίσεις για ταχεία βελτιστοποίηση, ενώ το LightGBM~\cite{ke2017lgbm} εισάγει
leaf-wise (αντί level-wise) ανάπτυξη δέντρων, επιτυγχάνοντας ταχύτερη εκπαίδευση με παρόμοια ή
καλύτερη ακρίβεια. Και οι δύο αλγόριθμοι περιλαμβάνουν ενσωματωμένη κανονικοποίηση (L1/L2 penalties)
που ελέγχει την υπεραρμογή.

Στη βιβλιογραφία EPF, το XGBoost επιτυγχάνει συνεχώς υψηλές επιδόσεις~\cite{lago2021,kanousis2025},
ενώ το LightGBM θεωρείται ελαφρώς βελτιωμένο για μεγάλα datasets~\cite{ke2017lgbm}.

\subsection{Support Vector Regression}

Η Support Vector Regression (SVR)~\cite{smola2004} εφαρμόζει τον αλγόριθμο SVM για παλινδρόμηση,
χρησιμοποιώντας $\varepsilon$-insensitive loss function. Η ισχύς της SVR έγκειται στη χρήση
kernel functions (RBF, πολυωνυμικό) που επιτρέπουν ισχυρή μη-γραμμική προσαρμογή. Ωστόσο, η SVR
κλιμακώνεται άσχημα με μεγάλα datasets ($O(n^2)$ training complexity), οπότε χρησιμοποιείται
συνήθως με υποδειγματοληψία.

\subsection{Νευρωνικά Δίκτυα --- MLP}

Τα Πολυεπίπεδα Perceptron (\textit{Multi-Layer Perceptrons}, MLP) αποτελούν τη βασική αρχιτεκτονική
βαθιάς μάθησης για πρόβλεψη χρονοσειρών~\cite{goodfellow2016deep}. Στο EPF εφαρμόζονται ως
feed-forward networks με διάφορα μεγέθη κρυφών στρωμάτων, activation functions (ReLU, tanh) και
regularization τεχνικές (dropout, weight decay). Σε σύγκριση με μοντέλα δέντρων, τα MLP τείνουν να
υπολείπονται χωρίς ενδελεχή ρύθμιση, αλλά μπορούν να αποδώσουν καλύτερα όταν οι σχέσεις είναι
πολύπλοκες και τα δεδομένα αρκετά.


\section{Βαθιά Μάθηση}
\label{sec:deep_learning}

\subsection{LSTM και RNN}

Τα Long Short-Term Memory (LSTM) δίκτυα~\cite{hochreiter1997lstm} είναι ειδικά σχεδιασμένα για
χρονοσειρές, αξιοποιώντας gates (input, forget, output) για επιλεκτική αποθήκευση μακροπρόθεσμων
εξαρτήσεων. Στο EPF, τα LSTM έχουν δείξει υποσχόμενα αποτελέσματα όταν το ορόσημο εισόδου (sequence
length) είναι αρκετά μεγάλο~\cite{kanousis2025}. Ο Kanousis (2025)~\cite{kanousis2025} χρησιμοποιεί
LSTM με 168-ωριαία ακολουθία εισόδου και 1 κρυφό στρώμα 24 κυττάρων με Dropout 20\%.

\subsection{Transformer-based Μοντέλα}

Τα τελευταία χρόνια, αρχιτεκτονικές βασισμένες σε Transformer~\cite{vaswani2017attention}, όπως το
Informer~\cite{zhou2021informer} και το PatchTST, εφαρμόζονται στο EPF. Αν και προσφέρουν
θεωρητικά πλεονεκτήματα (self-attention για μακρο-εξαρτήσεις), στην πράξη δεν ξεπερνούν πάντοτε τα
απλούστερα μοντέλα δέντρων σε βραχυπρόθεσμες ωριαίες προβλέψεις~\cite{zeng2023transformers}.


\section{Στρατηγικές Πολυ-βήματης Πρόβλεψης}
\label{sec:forecasting_strategies}

Ένα κρίσιμο ζήτημα που συχνά παραλείπεται στη βιβλιογραφία είναι η επιλογή της \textit{στρατηγικής
πρόβλεψης} --- δηλαδή του τρόπου με τον οποίο το μοντέλο παράγει προβλέψεις για ορίζοντα $H > 1$
βημάτων.

\subsection{Teacher-Forcing (Closed-Loop)}

Στη στρατηγική Teacher-Forcing (TF), η οποία ονομάζεται επίσης Closed-Loop ή Iterated One-Step
Ahead, το μοντέλο εκπαιδεύεται για πρόβλεψη ενός βήματος χρησιμοποιώντας \textit{πραγματικές}
καθυστερημένες τιμές ως χαρακτηριστικά εισόδου~\cite{bontempi2012strategies}. Κατά την αξιολόγηση,
εφόσον υπάρχουν πραγματικές τιμές (δηλαδή σε offline scenario), χρησιμοποιούνται πάλι οι
πραγματικές τιμές. Αυτό δίνει στρατηγικό πλεονέκτημα: δεν υπάρχει συσσώρευση σφάλματος.

Παρόλα αυτά, η TF στρατηγική δεν εφαρμόζεται σε πραγματικό χρόνο χωρίς πρόβλεψη λαγ τιμών (οπότε
πρέπει να συνδυαστεί με κάποιο oracle ή να περιοριστεί σε μοντέλα που δεν χρειάζονται μελλοντικές
τιμές ως χαρακτηριστικά). Στο πλαίσιο της παρούσας εργασίας, η TF αξιολογείται ως θεωρητικό άνω
όριο επίδοσης.

\subsection{Recursive / Open-Loop}

Στη Recursive στρατηγική (Rec) ή Open-Loop (OL), ένα μοντέλο πρόβλεψης ενός βήματος χρησιμοποιεί
τις \textit{δικές του προβλέψεις} ως καθυστερημένα χαρακτηριστικά για να παράξει την επόμενη
πρόβλεψη~\cite{taieb2012review}. Αυτό αντιστοιχεί στο αυτο-αναφορικό (autoregressive) σχήμα που
χρησιμοποιείται σε πραγματικό χρόνο. Το κύριο μειονέκτημα είναι η \textit{συσσώρευση σφάλματος}:
οι λανθασμένες προβλέψεις του παρελθόντος επηρεάζουν τις μελλοντικές.

Η τεχνική \textit{Scheduled Sampling} (SS)~\cite{bengio2015scheduled} αντιμετωπίζει αυτό
εκθέτοντας το μοντέλο κατά την εκπαίδευση σε συνδυασμό πραγματικών τιμών και προβλέψεων, με
πιθανότητα που μεταβαίνει σταδιακά από 100\% πραγματικές τιμές (TF) σε αποκλειστικά προβλέψεις (Rec).
Η τεχνική αυτή βελτιώνει την ανθεκτικότητα του μοντέλου σε θόρυβο εισόδου.

\subsection{MIMO}

Η στρατηγική MIMO (Multi-Input Multi-Output)~\cite{taieb2012review} εκπαιδεύει ένα μοντέλο που
εξάγει ταυτόχρονα $H$ τιμές για ολόκληρο τον ορίζοντα πρόβλεψης ως διάνυσμα εξόδου. Στο MIMO
δεν υπάρχει recursion: το μοντέλο "βλέπει" τις τελευταίες διαθέσιμες τιμές και παράγει απευθείας
το σύνολο $\{y_{t+1}, y_{t+2}, \ldots, y_{t+H}\}$. Αυτό αποτρέπει τη συσσώρευση σφάλματος αλλά
καθιστά την εργασία εγγενώς δυσκολότερη σε σύγκριση με TF, δεδομένου ότι ο ορίζοντας προβλέπεται
χωρίς ανανέωση πληροφορίας.

Η παραλλαγή \textit{Direct}~\cite{taieb2012review} εκπαιδεύει ξεχωριστό μοντέλο ανά βήμα ορίζοντα
($H$ ξεχωριστά μοντέλα), αποφεύγοντας την κοινή εξαγωγή αλλά διατηρώντας ανεξαρτησία ανά ορίζοντα.

\subsection{Walk-Forward Retraining}

Στη στρατηγική Walk-Forward Retraining (MR), η εκπαίδευση ανανεώνεται κυλιόμενα (rolling) σε τακτά
χρονικά διαστήματα (π.χ.\ μηνιαία), ώστε το μοντέλο να προσαρμόζεται σε μεταβαλλόμενες
συνθήκες~\cite{diebold2015forecasting}. Αυτό είναι ιδιαίτερα σημαντικό για αγορές που
χαρακτηρίζονται από \textit{distribution shifts} λόγω αλλαγής ενεργειακού μίγματος, τιμών
φυσικού αερίου κ.λπ.


\section{Βιβλιογραφία για την Ελληνική Αγορά}
\label{sec:greek_literature}

Η βιβλιογραφία για την Ελληνική αγορά παραμένει σχετικά περιορισμένη. Ο Kanousis (2025)~\cite{kanousis2025}
είναι από τις πρώτες εργασίες που εφαρμόζει σύγχρονους αλγόριθμους (RF, XGBoost, SARIMA, LSTM) για
την πρόβλεψη τιμής DAM της ελληνικής αγοράς για ολόκληρο το 2024, επιτυγχάνοντας XGBoost MAE ~6,5
€/MWh. Χρησιμοποιεί day-ahead προβλέψεις φορτίου, ηλιακής και αιολικής παραγωγής ως εξωγενείς
μεταβλητές. Ο Magklaras~\cite{magklaras2022} εξετάζει βραχυπρόθεσμη πρόβλεψη ηλεκτρικού φορτίου σε
γραμμές μέσης τάσης Ελλάδος και Ιταλίας, χρησιμοποιώντας LSTM, MLP, CNN και ANN.

Γενικότερα, για ευρωπαϊκές αγορές, η εργασία των Marcjasz et al.\ (2018)~\cite{marcjasz2018} δείχνει
ότι ML μοντέλα (ιδίως GBM) υπερτερούν σταθερά έναντι στατιστικών μοντέλων στην ωριαία πρόβλεψη
DAM τιμής. Ο Uniejewski et al.\ (2019)~\cite{uniejewski2019} υποστηρίζει ότι η κανονικοποίηση
εξόδου (variance stabilization) βελτιώνει σημαντικά τα αποτελέσματα.


\section{Σύγκριση με την Παρούσα Εργασία}
\label{sec:comparison}

\begin{table}[H]
\centering
\caption{Σύγκριση με εκπροσωπευτικές εργασίες της βιβλιογραφίας.}
\label{tab:lit_comparison}
\small
\begin{tabularx}{\linewidth}{p{2.8cm} X X X X X}
\toprule
\textbf{Εργασία} & \textbf{Αγορά} & \textbf{Εργασία} & \textbf{Αλγόριθμοι} & \textbf{Στρατηγικές} & \textbf{DM test} \\
\midrule
Lago et al.\ 2021 & Πολλές (EU) & Τιμή & Πολλοί & 1-step & Ναι \\
Kanousis 2025 & Ελλάδα & Τιμή & RF, XGB, SARIMA, LSTM & rolling TF & Ναι \\
Magklaras 2022 & Ελλάδα/Ιταλία & Φορτίο & LSTM, MLP, CNN & 1-step & Όχι \\
\midrule
\textbf{Παρούσα} & \textbf{Ελλάδα} & \textbf{Τιμή + Φορτίο} & \textbf{LGBM, XGB, RF, SVR, MLP} & \textbf{TF, Rec/OL, MIMO, MR} & \textbf{Ναι} \\
\bottomrule
\end{tabularx}
\end{table}

Η παρούσα εργασία διαφοροποιείται από τις υφιστάμενες κυρίως ως προς:
(α)~τη σύγκριση \textit{τεσσάρων} διαφορετικών στρατηγικών πρόβλεψης,
(β)~τη \textit{συνδυαστική} αξιολόγηση τιμής και φορτίου,
(γ)~την αξιολόγηση σε πολλαπλές χρονικές κλίμακες (1 εβδομάδα, 1 μήνας, 3 μήνες), και
(δ)~τη χρήση Monthly Retraining ως στρατηγική που εισάγει δυναμική προσαρμογή.
